% Generated draft: requires mathematical and source review.
\documentclass[11pt]{article}
\usepackage[margin=25mm]{geometry}
\usepackage{fontspec}
\setmainfont{Latin Modern Roman}
\usepackage{amsmath,amssymb,mathtools,unicode-math}
\setmathfont{Latin Modern Math}
\usepackage{enumitem,needspace,xurl,hyperref}
\setlength{\parindent}{0pt}
\setlength{\parskip}{5pt}
\setlength{\emergencystretch}{3em}
\newcommand{\N}{\mathbb{N}}
\newcommand{\Z}{\mathbb{Z}}
\newcommand{\Q}{\mathbb{Q}}
\newcommand{\R}{\mathbb{R}}
\newcommand{\C}{\mathbb{C}}
\newcommand{\F}{\mathbb{F}}
\newcommand{\A}{\mathbb{A}}
\newcommand{\PP}{\mathbb{P}}
\newcommand{\E}{\mathbb{E}}
\newcommand{\eps}{\varepsilon}
\newcommand{\dd}{\,\mathrm{d}}
\begin{document}
\section*{UnsolvedMath ID: 30001683}
\textbf{Source problem code:} catalogue entry
\subsection{Definitions and mathematical statement}
Let $F$ denote Thompson's group: the group, under composition, of increasing piecewise linear homeomorphisms of $[0,1]$ with finitely many breakpoints, all dyadic rational, and slopes in $\{2^k:k\in\Z\}$. Here a dyadic rational is a number $m/2^r$ with $m\in\Z$ and $r\in\Z_{\geq 0}$. Say that $F$ is amenable if there exists a finitely additive probability measure $\mu:\mathcal P(F)\to[0,1]$ such that $\mu(gE)=\mu(E)$ for every $g\in F$ and every $E\subseteq F$.

Let $\mathcal T$ be the set of finite formal terms generated by the constant symbol $1$ and the binary operation symbol $\star$: thus $1\in\mathcal T$, and $A\star B\in\mathcal T$ whenever $A,B\in\mathcal T$. For this draft, interpret being related by the displayed identity as the least congruence $\sim$ on $\mathcal T$ containing
\[
 A\star(B\star C)\sim (A\star B)\star(A\star C)
 \qquad(A,B,C\in\mathcal T).
\]
A congruence here is an equivalence relation preserved by $\star$ in both arguments. This specifies closure under replacements inside larger terms; the supplied record does not explicitly state that closure convention.

Write $\N$ for the natural-number domain in the supplied record, whose starting index is unspecified. For a strictly increasing nonidentity function $f:\N\to\N$, interpret its critical point as
\[
 \operatorname{crit}(f)=\min\{n\in\N:f(n)\ne n\}.
\]
This is the conventional first-moved-point interpretation; the supplied record does not define $\operatorname{crit}$, so this interpretation requires confirmation against the report.

With these qualifications, the proposed equivalence is
\[
 F\text{ is amenable}
 \quad\Longleftrightarrow\quad
 \text{there exists a family }(f_T)_{T\in\mathcal T}\text{ satisfying the conditions below}.
\]
For every $T\in\mathcal T$, the function $f_T:\N\to\N$ is strictly increasing and is not the identity. The remaining conditions are
\begin{align*}
 S\sim T&\ \Longrightarrow\ f_S=f_T
 &&(S,T\in\mathcal T),\\
 \operatorname{crit}(f_{A\star B})&=f_A\bigl(\operatorname{crit}(f_B)\bigr)
 &&(A,B\in\mathcal T),\\
 \{\operatorname{crit}(f_T):T\in\mathcal T\}&=\N.
\end{align*}
The last condition means that every natural number occurs as the critical point of at least one member of the family. The equivalence is posed as a question, not asserted as a theorem.
\subsection{Short English statement}
Does amenability of Thompson's group $F$ hold exactly when there is a family of strictly increasing nonidentity maps of $\N$, indexed by variable-free terms, respecting the displayed distributive identity and critical-point recursion, with every natural number occurring as a critical point? The meanings of critical point, identity-generated equivalence, and the starting index of $\N$ require source confirmation.
\subsection{Sources}
\begin{itemize}
\item Local UnsolvedMath registry: supplied complete record with permanent ID $30001683$, title ``Functional Characterization of Thompson Group Amenability,'' source problem code \texttt{catalogue entry}, in the Oberwolfach Reports Open Problems collection. The supplied proposed year is 2011; no proposer is named.
\item Supplied report citation: \emph{Set Theory} (2011), Oberwolfach Reports, pp.~85--140, DOI \texttt{10.4171/owr/2011/02}. \url{https://doi.org/10.4171/owr/2011/02}. This is also the supplied literature source labelled ``Oberwolfach problem statement.''
\item Supplied literature source labelled ``current amenability survey'': \url{https://arxiv.org/abs/2305.07113}. No author or title is supplied for this reference.
\item Literature assessment dated exactly \texttt{2026-08-21}: the supplied curation labels the problem open and describes the proposed equivalence as conjectural alongside the unresolved amenability problem for $F$. This is supplied dated curation, not an independently verified literature assessment. The background explicitly requests independent verification of time-sensitive status and cited claims.
\item The registry supplies the statement-status label \texttt{corrected\_verified} and reports that the official report was checked and the properties defining the function system were restored. These are supplied provenance claims; this draft does not independently verify them. The shorter original statement mentions only the equivalence and the term-indexed system; the detailed conditions above come from the supplied expanded statement.
\end{itemize}
\end{document}
